\part*{Bibliographie}
\addcontentsline{toc}{part}{Bibliographie}
\markboth{}{}

Les références ci-dessous rassemblent l'ensemble des sources citées en note au fil de l'ouvrage.

\noindent\hangindent=1.5em\hangafter=1 Appelbaum, S. H., Habashy, S., Malo, J. L., \& Shafiq, H. (2012). Back to the future: revisiting Kotter's 1996 change model. Journal of Management Development, 31(8), 764-782.\par\medskip
\noindent\hangindent=1.5em\hangafter=1 Bennis, W. (2009). On Becoming a Leader. Basic Books.\par\medskip
\noindent\hangindent=1.5em\hangafter=1 Berkun, S. (2013). The Year Without Pants: WordPress.com and the Future of Work. Jossey-Bass.\par\medskip
\noindent\hangindent=1.5em\hangafter=1 Bolden, R. (2011). Distributed leadership in organizations: A review of theory and research. International Journal of Management Reviews, 13(3), 251-269.\par\medskip
\noindent\hangindent=1.5em\hangafter=1 Bourdieu, P., \& Passeron, J.-C. (1970). La Reproduction. Éléments pour une théorie du système d'enseignement. Les Éditions de Minuit.\par\medskip
\noindent\hangindent=1.5em\hangafter=1 Bouville, G. (2008). La résistance au changement: réflexion théorique et expérience vécue. Revue française de gestion, (181), 111-129.\par\medskip
\noindent\hangindent=1.5em\hangafter=1 Cappelli, P., \& Tavis, A. (2016). The Performance Management Revolution. Harvard Business Review, 94(10), 58-67.\par\medskip
\noindent\hangindent=1.5em\hangafter=1 Carney, B. M., \& Getz, I. (2009). Freedom, Inc.: Free Your Employees and Let Them Lead Your Business for Greater Success. Crown Business.\par\medskip
\noindent\hangindent=1.5em\hangafter=1 Crozier, M., \& Friedberg, E. (1977). L'acteur et le système: Les contraintes de l'action collective. Seuil.\par\medskip
\noindent\hangindent=1.5em\hangafter=1 D'Iribarne, P. (1989). La logique de l'honneur: Gestion des entreprises et traditions nationales. Seuil.\par\medskip
\noindent\hangindent=1.5em\hangafter=1 Denning, S. (2018). The Age of Agile. Amacom.\par\medskip
\noindent\hangindent=1.5em\hangafter=1 Drucker, P. (1954). The Practice of Management. Harper \& Brothers.\par\medskip
\noindent\hangindent=1.5em\hangafter=1 Dupuy, F. (2015). La faillite de la pensée managériale: Lost in management 2. Seuil.\par\medskip
\noindent\hangindent=1.5em\hangafter=1 Endenburg, G. (1998). Sociocracy: The organization of decision-making, "no objection!" as the principle of sociocracy. Eburon.\par\medskip
\noindent\hangindent=1.5em\hangafter=1 Fleming, P. (2019). The Worst Is Yet to Come: A Post-Capitalist Survival Guide. Repeater Books.\par\medskip
\noindent\hangindent=1.5em\hangafter=1 Foss, N. J., \& Klein, P. G. (2014). Why Managers Still Matter. MIT Sloan Management Review, 56(1), 73-80.\par\medskip
\noindent\hangindent=1.5em\hangafter=1 Gandini, A. (2019). Labour process theory and the gig economy. Human Relations, 72(6), 1039–1056. https://doi.org/10.1177/0018726718792867\par\medskip
\noindent\hangindent=1.5em\hangafter=1 Gérard, A. (2017). Le patron qui ne voulait plus être chef. Flammarion.\par\medskip
\noindent\hangindent=1.5em\hangafter=1 Getz, I. (2009). Liberating Leadership: How the Initiative-Freeing Radical Organizational Form Has Been Successfully Adopted. California Management Review, 51(4), 32-58.\par\medskip
\noindent\hangindent=1.5em\hangafter=1 Getz, I. (2017). L'entreprise libérée. Fayard.\par\medskip
\noindent\hangindent=1.5em\hangafter=1 Goulston, M. (2009). Just Listen: Discover the Secret to Getting Through to Absolutely Anyone. Amacom.\par\medskip
\noindent\hangindent=1.5em\hangafter=1 Greenleaf, R. K. (1977). Servant leadership: A journey into the nature of legitimate power and greatness. Paulist Press.\par\medskip
\noindent\hangindent=1.5em\hangafter=1 Hamel, G. (2007). The Future of Management. Harvard Business School Press.\par\medskip
\noindent\hangindent=1.5em\hangafter=1 Hinds, P. J., \& Kiesler, S. (Eds.). (2002). Distributed work. MIT Press.\par\medskip
\noindent\hangindent=1.5em\hangafter=1 Hofstede, G. (2001). Culture's Consequences: Comparing Values, Behaviors, Institutions and Organizations Across Nations. Sage.\par\medskip
\noindent\hangindent=1.5em\hangafter=1 Hollander, E. P. (1992). Leadership, followership, self, and others. Leadership Quarterly, 3(1), 43-54.\par\medskip
\noindent\hangindent=1.5em\hangafter=1 Hsieh, T. (2010). Delivering happiness: A path to profits, passion, and purpose. Hachette Books.\par\medskip
\noindent\hangindent=1.5em\hangafter=1 Huy, Q. N. (2001). In praise of middle managers. Harvard Business Review, 79(8), 72-79.\par\medskip
\noindent\hangindent=1.5em\hangafter=1 Kotter, J. P. (1996). Leading Change. Harvard Business School Press.\par\medskip
\noindent\hangindent=1.5em\hangafter=1 Laloux, F. (2014). Reinventing Organizations: A Guide to Creating Organizations Inspired by the Next Stage of Human Consciousness. Nelson Parker.\par\medskip
\noindent\hangindent=1.5em\hangafter=1 Langfred, C. W., \& Moye, N. A. (2004). Effects of task autonomy on performance: an extended model considering motivational, informational, and structural mechanisms. Journal of Applied Psychology, 89(6), 934-945.\par\medskip
\noindent\hangindent=1.5em\hangafter=1 Lee, M. Y., \& Edmondson, A. C. (2017). Self-managing organizations: Exploring the limits of less-hierarchical organizing. Research in Organizational Behavior, 37, 35-58.\par\medskip
\noindent\hangindent=1.5em\hangafter=1 Mintzberg, H. (2004). Managers Not MBAs: A Hard Look at the Soft Practice of Managing and Management Development. Berrett-Koehler Publishers.\par\medskip
\noindent\hangindent=1.5em\hangafter=1 Mullenweg, M. (2020). The Five Levels of Autonomy. Ma.tt.\par\medskip
\noindent\hangindent=1.5em\hangafter=1 OCDE (2019). OECD Employment Outlook 2019: The Future of Work. OECD Publishing.\par\medskip
\noindent\hangindent=1.5em\hangafter=1 Puranam, P., \& Håkonsson, D. D. (2015). Valve's way. Journal of Organization Design, 4(2), 2-4.\par\medskip
\noindent\hangindent=1.5em\hangafter=1 Robertson, B. J. (2015). Holacracy: The New Management System for a Rapidly Changing World. Henry Holt and Co.\par\medskip
\noindent\hangindent=1.5em\hangafter=1 Rothschild, J., \& Whitt, J. A. (1986). The cooperative workplace. Cambridge University Press.\par\medskip
\noindent\hangindent=1.5em\hangafter=1 Sartre, J.-P. (1946). L'existentialisme est un humanisme. Éditions Nagel.\par\medskip
\noindent\hangindent=1.5em\hangafter=1 Semler, R. (2003). The Seven-Day Weekend: Changing the Way Work Works. Portfolio.\par\medskip
\noindent\hangindent=1.5em\hangafter=1 Stø, T., \& Ims, K. J. (2018). Freedom Inc. revisited: creating a freedom-based and healthier organization. Society and Business Review.\par\medskip
\noindent\hangindent=1.5em\hangafter=1 Teerlink, R., \& Ozley, L. (2000). More Than a Motorcycle: The Leadership Journey at Harley-Davidson. Harvard Business School Press.\par\medskip
\noindent\hangindent=1.5em\hangafter=1 Touraine, A. (1978). La voix et le regard. Seuil.\par\medskip
\noindent\hangindent=1.5em\hangafter=1 Valve Corporation (2012). Handbook for New Employees.\par\medskip
\noindent\hangindent=1.5em\hangafter=1 Zobrist, J.-F. (2012). La belle histoire de FAVI: L'entreprise qui croit que l'homme est bon. Humanisme \& Organisations.\par\medskip
